\subsection{赋值除环上的拓扑向量空间中原点的邻域}

定义 3. --- 设 K 是赋值除环，E 是 K 上的左向量空间；我们称 E 的子集 M 是均衡的，如果对所有 \(x \in M\) 及所有 \(\lambda \in K\)，只要 \(|\lambda| \leqslant 1\)，就有 \(\lambda x \in M\)（换言之，如果 \(\lambda M \subset M\) 只要 \(|\lambda| \leqslant 1\)）。

命题 2. --- 在赋值除环 K 上的拓扑向量空间 E 中，均衡集 M 的闭包是均衡集。

若 B 是由 \(\xi\in K\)（满足 \(|\xi|\leqslant1\)）组成的集，则 B 在 K 中闭。而 \(B\times M\) 被连续映射 \((\lambda,x)\mapsto\lambda x\) 映入 M；因此 \(B\times\overline{M}\) 被映入 \(\overline{M}\) (GT, I, §2.1, th. 1)，这就证明 \(\overline{M}\) 是均衡的。

当 M 是赋值除环 K 上的向量空间 E 中的任意集合时，集合 \(M_{1}\)（即诸 \(\lambda x\)，其中 \(x \in M\)、\(\lambda \in K\) 且 \(|\lambda| \leqslant 1\)）显然是包含 M 的最小均衡集；\(M_{1}\) 称为 M 的均衡包。

命题 3. --- 设 K 是局部紧的非离散赋值除环，E 是 K 上的 Hausdorff 拓扑向量空间（相应地，拓扑向量空间）。对 E 中每个紧（相应地，预紧）集 H，H 的均衡包是紧的（相应地，预紧的）。

若 B 表示 K 中的球 \(|\xi| \leqslant 1\)，则 H 的均衡包就是 \(H_{1}\)，即 \(B \times H\) 在连续映射 \(m: (\lambda, x) \mapsto \lambda x\) 下的像。若 E 是 Hausdorff 的，B 与 H 都是紧的，则 \(B \times H\) 从而 \(H_{1}\) 也是紧的。若 H 是预紧的，则 m 在 \(B \times H\) 上的限制是一致连续的 (I, 第 5 页, 注记 2)，而由于 \(B \times H\) 预紧，其在 m 下的像也是预紧的 (GT, II, § 4.2, prop. 2)。

注意，闭集的均衡包未必闭。例如在 \(\mathbf{R}^2\) 中，由方程 \(xy = 1\) 定义的双曲线的均衡包不是闭的。

E 中一族均衡集的并是均衡的，由此推出，对 E 的每个集合 M，存在包含在 M 中的最大均衡子集 N，称为 M 的均衡核；而且 N 非空当且仅当 \(0 \in M\)。说 \(x \in N\) 意味着对所有 \(\lambda \in K\)，只要 \(|\lambda| \leqslant 1\)，就有 \(\lambda x \in M\)，或者（当 \(0 \in M\) 时）等价地，对所有 \(\mu \in K\)，只要 \(|\mu| \geqslant 1\)，就有 \(x \in \mu M\)。若 \(0 \in M\)，则 M 的均衡核 N 就是交集 \(\bigcap_{|\mu| \geqslant 1} \mu M\)。这特别表明，若 M 闭，则 N 也闭。

定义 4. --- 设 K 是非离散赋值除环，E 是 K 上的左向量空间且带有两个子集 A 和 B。我们称 A 吸收 B，如果存在 \(\alpha > 0\)，使得 \(\lambda A \supset B\) 对每个 \(\lambda \in K\) 只要 \(|\lambda| \geqslant \alpha\) 就成立（等价地，如果 \(\mu B \subset A\) 对 \(\mu \neq 0\) 且 \(|\mu| \leqslant \alpha^{-1}\) 成立）。E 的集合 A 称为吸收集，如果它吸收每个由单独一点组成的集合。

设 A 是 E 的均衡集；为使 A 吸收 E 的集合 B，只需存在 \(\lambda \neq 0\) 使 \(\lambda A \supset B\)；事实上，当 \(|\mu| \geqslant |\lambda|\) 时，有 \(\lambda A = (\lambda \mu^{-1}) \mu A\)，而由于 \(\mu A\) 均衡且 \(|\lambda \mu^{-1}| \leqslant 1\)，由此推出 \(\lambda A \subset \mu A\)，从而 \(B \subset \mu A\)。特别地，为使 E 的均衡集 A 是吸收集，必须且只需：对每个 \(x \in E\)，存在 K 中的 \(\lambda \neq 0\) 使 \(\lambda x \in A\)。E 的每个吸收集生成向量空间 E。吸收集的每个有限交是吸收集。

命题 4. --- 在非离散赋值除环 K 上的拓扑向量空间 E 中，存在 0 的闭邻域组成的一个基本邻域系 \(\mathfrak{B}\)，使得：

\((\mathrm{EV}_{\mathrm{I}})\) 每个 \(\mathbf{V}\in \mathfrak{B}\) 都是均衡的且是吸收的。

(EV \(_{II}\) ) 对每个 V \(\in\) B 及 K 中的 \(\lambda \neq 0\)，有 \(\lambda\) V \(\in\) B（B 在比不为零的位似下不变）。

\((\mathrm{EV}_{\mathrm{III}})\) 对每个 \(\mathbf{V}\in \mathfrak{B}\)，存在 \(\mathbf{W}\in \mathfrak{B}\) 使得 \(\mathbf{W} + \mathbf{W}\subset \mathbf{V}\)。

反之，设 E 是 K 上的向量空间，V 是 E 上满足条件 \((\mathrm{EV}_{\mathrm{I}})\)、\((\mathrm{EV}_{\mathrm{II}})\) 与 \((\mathrm{EV}_{\mathrm{III}})\) 的一个滤子基。则在 E 上存在（且唯一）一个与 E 的向量空间结构相容的拓扑，使 V 构成 0 的一个基本邻域系。

由公理 \((\mathrm{EVT}_{\mathrm{III}}^{\prime})\)，我们首先证明：0 的邻域 V 的均衡核 \(V_{1}\) 本身是 0 的邻域。因为存在 \(\alpha > 0\) 及 0 的一个邻域 W，使得当 \(|\lambda| \leqslant \alpha\) 且 \(x \in W\) 时，有 \(\lambda x \in V\)。由于 K 非离散，存在 K 中的 \(\mu \neq 0\) 满足 \(|\mu| \leqslant \alpha\)，而 \(\mu W\) 是 0 的一个邻域且 \(\mu W \subset V\)。又若 \(v \in K\) 且 \(|v| \leqslant 1\)，则 \(|v\mu| \leqslant \alpha\)，从而 \(v\mu W \supset V\)。于是 \(\mu W \supset V_{1}\)，故 \(V_{1}\) 是 0 的邻域。又由于 V 闭，\(V_{1}\) 也闭。因此，0 的闭均衡邻域所成的集 B 构成 E 中 0 的一个基本邻域系。由 \((\mathrm{EVT}_{\mathrm{I}}^{\prime})\)，0 的每个邻域都是吸收的；此外 B 满足 \((\mathrm{EV}_{\mathrm{II}})\) (参见 I, 第 3 页, cor.)；最后，由于 \((x, y) \mapsto x + y\) 在点 \((0, 0)\) 处连续，E 中 0 的每个基本邻域系都满足 \((\mathrm{EV}_{\mathrm{III}})\)。于是集 B 满足命题中的条件。

现在设 E 是 K 上的向量空间，V 是 E 上满足 \((\mathrm{EV}_{\mathrm{I}})\)、\((\mathrm{EV}_{\mathrm{II}})\) 与 \((\mathrm{EV}_{\mathrm{III}})\) 的一个滤子基。公理 \((\mathrm{EV}_{\mathrm{I}})\) 首先表明，对所有 \(V \in V\)，有 -V = V 及 \(0 \in V\)；这些关系式与公理 \((\mathrm{EV}_{\mathrm{III}})\) 表明，V 构成 E 上与 E 的加法群结构相容的某个拓扑 (GT, III, § 1.2) 下 0 的基本邻域系。另一方面，公理 (EVT \(_{I}\) )、(EVT \(_{II}\) ) 与 (EVT \(_{III}\) ) 是 (EV \(_{I}\) ) 与 (EV \(_{II}\) ) 的直接推论，因此上面定义的拓扑满足公理 (EVT)，命题得证。

注记. --- 1) 在非离散赋值除环上的赋范空间中，以 0 为中心的全体开球（相应地，闭球）构成 0 的一个基本邻域系，它满足条件 \((\mathrm{EV}_{\mathrm{I}})\)、\((\mathrm{EV}_{\mathrm{II}})\) 与 \((\mathrm{EV}_{\mathrm{III}})\)。

\begin{enumerate}
\def\labelenumi{\arabic{enumi})}
\setcounter{enumi}{1}
\tightlist
\item
  当标量除环 \(\mathbf{K}\) 是 \(\mathbf{R}\) 域或 \(\mathbf{C}\) 域时，每个滤子基 \(\mathfrak{B}\)（在 \(\mathbf{E}\) 上）若只满足两条公理 \((\mathrm{EV}_{\mathbf{I}})\) 与 \((\mathrm{EV}_{\mathrm{III}})\)，则它是某个与 \(\mathbf{E}\) 的向量空间结构相容的拓扑下 0 的基本邻域系。事实上，只需证明：在这些条件下，对每个 \(\lambda \neq 0\)（在 \(\mathbf{K}\) 中）及每个 \(\mathrm{V} \in \mathfrak{B}\)，存在 \(\mathrm{W} \in \mathfrak{B}\) 使得 \(\lambda \mathrm{W} \subset \mathrm{V}\)。由 \((\mathrm{EV}_{\mathrm{III}})\)，存在 \(\mathrm{W}_{1} \in \mathfrak{B}\) 满足 \(2\mathrm{W}_{1} \subset \mathrm{V}\)，并且我们用归纳法推出，对每个正整数 \(n\)，存在 \(\mathrm{W}_{n} \in \mathfrak{B}\) 使得 \(2^{n}\mathrm{W}_{n} \subset \mathrm{V}\)。由于 \(\mathrm{V}\) 是均衡的，只要取 \(n\) 大到使 \(2^{n} = |2^{n}| > |\lambda|\)，则 \(\mathrm{W} = \mathrm{W}_{n}\) 就满足条件，此即所需。
\end{enumerate}

这个结果并非对每个非离散赋值除环 K 都成立，因为在这种除环中，对每个正整数 m 不再必有 \(|me| = m\)（\(e\) 表示除环的单位元；参见 I, 第 22 页, 习题 1）。

\begin{enumerate}
\def\labelenumi{\arabic{enumi})}
\setcounter{enumi}{2}
\tightlist
\item
  若 K 是离散除环，则条件 \((\mathrm{EVT}_{\mathrm{I}}^{\prime})\) 与 \((\mathrm{EVT}_{\mathrm{III}}^{\prime})\) 对 E 上的任何拓扑都成立。按命题 4 的方式论证，容易看出：若 E 是 K 上的拓扑向量空间，则存在 E 中 0 的闭邻域组成的一个基本邻域系 B，满足条件 \((\mathrm{EV}_{\mathrm{II}})\) 与 \((\mathrm{EV}_{\mathrm{III}})\)。反之，若 K 上向量空间 E 上的滤子基 B 使得 0 属于 B 的所有集合，且 \((\mathrm{EV}_{\mathrm{II}})\)、\((\mathrm{EV}_{\mathrm{III}})\) 成立，则 B 是与 E 的向量空间结构相容的某个拓扑下 0 的基本邻域系。
\end{enumerate}

